\section{Therapist, delivery and real-world factors}\label{sec:delivery}
\paragraph{Therapists} Therapist effects are real but modest: about 5\% of outcome variance on average (range 0.2--29\%; \citealp{johns2019therapist}), 6--7\% among low-intensity workers in IAPT \citep{firth2015pwp}, and in one large routine dataset individual therapists' recovery rates ranged from 23.5\% to 95.6\% \citep{saxon2012therapist}. Whether delivering CBT competently or adherently matters is unsettled: a meta-analysis found adherence--outcome and competence--outcome correlations near zero ($r=.02$ and $.07$; modalities pooled; \citealp{webb2010adherence}), while single CBT studies find competence effects mostly at the extremes or in subgroups \citep{branson2015competence,strunk2010competence,kuyken2008competence}. Experience alone did not improve outcomes across 6,591 patients and 170 therapists \citep{goldberg2016experience}. The alliance--outcome association is $r=.278$ across 295 studies and the same for internet therapy \citep{fluckiger2018alliance}, and sits at the therapist rather than the patient level \citep{baldwin2007alliance}; it is correlational and not CBT-specific.

\paragraph{Drift and training} Reviews argue that clinicians under-use core behavioural components and drift from ``doing'' to ``talking'' therapy \citep{waller2009drift,waller2016driftredux}; a survey found most licensed psychologists did not use exposure for PTSD \citep{becker2004exposure} (dated, US). We found no pooled prevalence of drift. Training changes behaviour only when it is active and supervised: unsupervised internet training did not change competence in a randomized trial \citep{rakovshik2016supervision}, in line with broader reviews \citep{beidas2010training,shafran2009gap}.

\paragraph{Dose, homework and schedule} Homework compliance correlates with outcome ($r=.26$, 23 studies; \citealp{mausbach2010homework}; $g=0.79$ for quantity, \citealp{kazantzis2016quality}), but patients who improve may do more homework. Improvement is curvilinear with dose, and most responders in routine care respond within 7 (low-intensity) or 14 (high-intensity) sessions \citep{robinson2020dose}; session frequency may matter more than total sessions \citep{cuijpers2013sessions}; intensified exposure scheduling gave equivalent outcomes with 32\% faster response \citep{pittig2021intensified}.

\paragraph{Does CBT survive contact with routine care?} Largely yes. Effectiveness studies for anxiety lie within the range of efficacy trials \citep{Stewart2009Anxiety}, those for depression are somewhat lower with about 25\% dropout \citep{Hans2013Depression}, and clinical representativeness did not erase effects \citep{shadish2000representative}. IAPT shows pre--post $d$ of 0.87 (depression) and 0.88 (anxiety) \citep{wakefield2021iapt}, though year-one reliable recovery was 40.3\% \citep{Gyani2013IAPT}, and a small independent diagnostic follow-up found only 9.2\% fully recovered (n = 90; single author; contested; \citealp{scott2018iapt}). Youth evidence-based practices beat usual care only modestly (mean effect 0.29; \citealp{weisz2013youthusual}).

\paragraph{Format and personalisation} Guided internet CBT is equivalent to face-to-face CBT ($g=0.05$; \citealp{carlbring2018icbt}), works in routine care \citep{etzelmueller2020routine}, and guided beats unguided only slightly at post-treatment \citep{karyotaki2021icbtnma}; unguided CBT has a small effect ($g=0.27$; \citealp{karyotaki2017selfguided}). A large pragmatic trial of supported computerised CBT in primary care found no benefit over usual care, with low uptake \citep{gilbody2015reeact}. Group and individual formats are equivalent when treatments are identical \citep{burlingame2016group}. Manualisation has no direct evidence of superiority \citep{truijens2019manual}; algorithm-matched care modestly beat stepped care (OR 1.40; \citealp{delgadillo2022stratified}; see also \citealp{firth2015stepped}).

\paragraph{Reading} R9 is partly supported. CBT works in routine care with lower and more variable yield, and the modifiable factors (supervision, dose, schedule, feedback, matching) are known, but causal evidence that poor delivery explains non-response is missing.
